\chapter{Analysis of braking events}
\label{ch:braking}


\section{Events}

A deceleration event is one continuous stretch of the deceleration state; the training window
yields about 126{,}700, of which 99.9\% are genuine decelerations once the held-brake
contamination was removed (Section~\ref{sec:heldbrake}). The mean event lasts 23.4\,s and the
median 20\,s (95th percentile 30\,s): mostly ordinary station stops, with a thin tail of longer
ones.

\section{Classes or continuum?}
\label{sec:continuum}

The original plan anticipated discrete braking classes: light, hard, emergency. The data was
allowed to decide instead. Clustering the intensity features and reading BIC, silhouette, and
resampling stability together turns up no natural class count: the BIC keeps improving as
clusters are added rather than settling at a minimum, and the silhouette never singles out a $k$
(Figure~\ref{fig:continuum}). Stability does rise again at $k=5$, which looks stable at a glance
but survives only by peeling off a cluster holding a few hundredths of a percent of the events,
unsupported by BIC or silhouette: the bump is an artefact of one tiny outlier group, not a regime.
The intensity histogram seems to pull the other way, with a dominant mild mode and a broader hump
at firmer braking (Figure~\ref{fig:intensity}), so on its own the density looks bimodal. But the
humps overlap without a clean gap, and the resampling scrambles any boundary drawn between them,
so the clustering never confirms the visual split. That tension is left as an open question; on
the available evidence the report treats intensity as a continuum, varying by degree rather than
by kind. No light/hard/emergency labels are introduced, then. Any such split is an \emph{operational
threshold} to be set with domain input, not a structure the data volunteers.

\begin{figure}[htbp]
  \centering
  \includegraphics[width=\linewidth]{cluster_model_selection.png}
  \caption{Model selection for braking intensity. BIC keeps improving with more groups and
  the silhouette falls away, so there is no natural number of intensity classes. The
  resampling stability rises again at $k=5$, but that solution only splits off a negligibly
  small cluster and no other criterion supports it, so the bump carries no interpretable
  meaning: intensity is a continuum.}
  \label{fig:continuum}
\end{figure}

\begin{figure}[htbp]
  \centering
  \includegraphics[width=1\linewidth]{intensity_distribution.png}
  \caption{Distribution of deceleration intensity. A dominant mild mode and a broader
  firmer-braking hump overlap without a clean gap; the cluster tests find no stable boundary
  between them, so intensity is treated as a continuum rather than as discrete classes.}
  \label{fig:intensity}
\end{figure}

\section{A deceleration-energy axis}
\label{sec:energy}

To compare stops on a physical scale, a specific deceleration energy is derived,
$\tfrac{1}{2}(v_0^2 - v_1^2)$, the kinetic energy dissipated per stop per unit mass. Being
kinematic, it carries the role \emph{velocity} and serves only as a target, never a predictor.
Adding it to the intensity features does not create classes (the continuum of
Section~\ref{sec:continuum} is unchanged) but gives a common wear axis, putting every stop on the
same energy-dissipated scale. The air-system sensors explain about 85\% of a stop's specific
energy out-of-sample (full tier, random forest), so the pneumatic record tracks dissipated energy
closely even without seeing the speed.

\section{Intensity regression}
\label{sec:regression}

The central cross-modal test predicts deceleration \emph{magnitude} from the non-speed sensors on
held-out events, against a guess-the-average baseline ($R^2 \approx 0$ by construction). A random forest on the full tier reaches $R^2 \approx 0.43$ for peak and $0.57$ for mean
deceleration; a linear model manages only 0.24 and 0.20, so most of the signal is non-linear
structure a straight fit misses. These are moderate results: the pneumatic sensors carry real but
partial information about how hard the train braked, not a near-perfect reconstruction. The scatter of predicted against actual intensity (Figure~\ref{fig:regscatter}) is positively
associated but sits in a narrower band than the targets: the gentlest stops are pushed up and the
hardest pulled down, flattening the cloud against the diagonal. This is regression toward the
average, where a squared-error model retreats when the inputs leave real variance unexplained, and
the per-event pneumatic signal is that noisy. Two other explanations were checked and rejected. Not sensor saturation: the actuation pressures
stay well inside range on hard stops, with only one to two percent of events near a ceiling. Not
the ten-second windowing either: adding event duration to the predictors moves $R^2$ by less than
a percentage point. The model did matter. A gradient-boosted tree scores 0.47 on mean deceleration
in a tighter band; the random forest reaches 0.57 and widens the spread by about a fifth, since
deep independent trees reach more extreme conditional means than a shallow boosted ensemble. Peak
deceleration holds at 0.43 either way, so its residual band is the limit of what these sensors
carry, not a model shortcoming. Both panels of Figure~\ref{fig:regscatter} use the random forest. A Lasso fit with its penalty cross-validated on the training events keeps 19 of the 22 air-system
features ($\alpha \approx 4\times10^{-6}$), a mix of actuation and auxiliary channels spread across
subsystems rather than concentrated in the brake command. The near-minimal penalty and broad
selection say the intensity signal is distributed, not localised, matching the auxiliary-only
classification result of Chapter~\ref{ch:classification}.

\begin{figure}[htbp]
  \centering
  \includegraphics[width=\linewidth]{decel_regression_scatter.png}
  \caption{Predicted versus actual deceleration intensity on held-out events (full tier,
  random forest). The association is positive, but predictions sit in a narrower band than
  the targets: the gentlest stops are over-predicted and the hardest under-predicted, the
  regression-toward-the-average that remains once the pneumatic sensors have given up their
  partial signal.}
  \label{fig:regscatter}
\end{figure}
